% v4 Methods 3.5: protocol, phases, metrics, uncertainty accounting
% (MC-8, MC-11). Constants from outputs/session35/mc_provenance.md.

\subsection{Protocol, phases and uncertainty accounting}
\label{sec:methods_protocol}

\rb{K}{S35-01}{Defines the three phases and every load metric used in the paper. Trim internally: the SSIM convention is already given in S2-06.}Every encounter is scored in three phases anchored to its impact frame:
lead-in $[t_{\mathrm{imp}} - 8,\, t_{\mathrm{imp}})$, impact
$[t_{\mathrm{imp}},\, t_{\mathrm{imp}} + 16)$ and relaxation
$[t_{\mathrm{imp}} + 16,\, t_{\mathrm{imp}} + 48)$, all half-open in
frames and clamped to the encounter. Load metrics are reported in three
complementary forms, because each corrects a failure mode of the others. The
coefficient of determination $R^2 = 1 - \mathrm{SSE}/\mathrm{SST}$ is pooled
over the phase frames, where $\mathrm{SSE}$ is the sum of squared errors of the
estimate and $\mathrm{SST}$ the sum of squared deviations of the truth about its
own mean, so $R^2$ measures skill against predicting that mean and is negative
whenever the estimate is worse than the mean would have been. It is scale free,
which is what makes it comparable across families, and that is also its failure
mode: at a phase where the truth varies a great deal between encounters, a large
$\mathrm{SST}$ flatters it. So we report alongside it the root-mean-square and
mean absolute errors in lift-coefficient units, which are not scale free and
expose exactly that case, and peak-relative quantities, which ask about the one
instant an extreme-gust application cares about. The peak diagnostics locate the
extremum of $|C_L|$ within the impact phase in the estimate and in the truth
separately, and report the value difference, its magnitude as a percentage of
the true peak, and the timing difference in convective units. Peak-region
readability uses a pooled $R^2$ over windows of half-width $8$ frames centred on
each encounter's true peak, against a persistence floor that holds the mean of
frames $[0, 25)$ constant; phase lag is the argmax of the normalised
cross-correlation within $\pm 20$ frames with parabolic sub-frame refinement,
positive when the estimate trails the truth. Decoded fields are scored by the
structural similarity index of \citet{wang2004ssim} on three masks, the
feathered near-body band of \S\ref{sec:methods_chang}, a wake window spanning $0
< x/c < 4.5$ and $|y/c| < 1.25$, and the full frame; its constants and its
dataset-dependent range are given in \S\ref{sec:flow_observables}. Linear-probe $R^2$ is
used throughout as a readability metric and never as an information metric: a
low value means the quantity is not linearly accessible in that latent, not that
it is absent. The control that motivates the distinction, a multilayer probe
that equalises what a linear read-out cannot reach, is reported alongside every
dimension comparison (\S\ref{sec:res_reconstruction}). An observable read at one
instant rather than over a window uses the same probe, fit on the training
window and evaluated on held-out encounters at the pre-impact, impact and
peak-lift frames ($\pm 2$), and is reported both as $R^2$ against the mean
observable at that instant and as a root-mean-square deviation in the
observable's own units, for the reason given above.\re

\rb{K}{S35-02}{Enumerates the six uncertainty sources and points at the one authoritative accounting. The sign-convention clause is still an open REVIEW-CLAIM.}Uncertainty enters the reported numbers from six places. Two are sampling: a
bootstrap over held-out encounters, and a clustered bootstrap over whole cases
for every paired comparison between families, since encounters of one case are
not independent of each other and treating them as such would understate the
interval. Three are training stochasticity: the encoder seed, the seed of an
operator fitted on a frozen representation, and the member-perturbation seed of
the filter, which are separated because they answer different questions, namely
whether a representation is reproducible, whether a downstream operator is, and
whether the estimator is. The sixth is the read-out, cross-validated over cases,
which bounds how much a probe result depends on the particular fit rather than
on the latent. Table~\ref{tab:uncertainty} gives the resampling or training unit
for each, its count, and which results carry it; that table is the authoritative
accounting, and it is there because no single number carries all six and which
apply depends on what was varied to produce it. Every stochastic number in the
paper states its count. Single-seed cells appear only in the dimension grid and
say so in their captions. The acceptance gates for the conditioning study and
for this campaign's follow-up runs were written and committed before any
corresponding result was read, and both declarations ship with the data record.
% REVIEW-CLAIM: (Stage 4) the following sentence states the archive
% (inventory) sign convention for the gust axis, which tenses against the
% rewrite rule "physical G everywhere; archive identifiers only in the data
% appendix". Most reporting is |G|-symmetric; the signed statements (the
% envelope asymmetry) need a one-time audit at the Stage 5 figure pass
% before the convention sentence is rewritten. Carlos to confirm.
The gust-intensity axis is
reported in the inventory sign convention, and no test-C quantity is used
for any selection anywhere in the pipeline.\re
